\section*{第 \ref{ch:g5:healthy-choices} 章——健康的选择}

\begin{solution}{exo:g5:healthy-choices:1}
\omterm{prop:g1:healthy-habits:needs}{多样的食物}和水：为每个\omterm{def:g5:first-look-at-cells:cell}{细胞}提供燃料
和建造材料，由\omterm{prop:g4:journey-of-food:absorb}{血液}输送。运动：\omterm{def:g3:bones-and-muscles:muscle}{肌肉}、骨、\omterm{def:g4:heart-and-blood:heart}{心脏}
和肺，越用越强健。\omterm{prop:g1:healthy-habits:needs}{睡眠}：\omterm{def:g5:brain-in-command:system}{大脑}的整理归档，
以及身体的修复。清洁——空气、双手、
牙齿：身体抵御\omterm{def:g1:healthy-habits:germs}{病菌}和污垢的防线。
\end{solution}

\begin{solution}{exo:g5:healthy-choices:2}
约十个小时。\omterm{def:g5:brain-in-command:system}{大脑}整理白天的经历（并做梦）；
身体修复并生长。
\end{solution}

\begin{solution}{exo:g5:healthy-choices:3}
一天的时间是固定的，因此看屏幕的时间只能从其他活动中挤出；
最常被挤掉的是户外运动和\omterm{prop:g1:healthy-habits:needs}{睡眠}。
\end{solution}

\begin{solution}{exo:g5:healthy-choices:4}
因为明亮的屏幕会告诉\omterm{def:g5:brain-in-command:system}{大脑}“仍是白天”，
使入睡推迟到那一小时之后——晚上看屏幕本身要占去一个小时，
还会占去夜晚的一部分。
\end{solution}

\begin{solution}{exo:g5:healthy-choices:5}
烟草烟雾先伤害肺部脆弱的气囊，
再加重\omterm{def:g4:heart-and-blood:heart}{心脏}的负担。酒精先作用于\omterm{def:g5:brain-in-command:system}{大脑}——\omterm{def:g5:brain-in-command:system}{大脑}是指挥中心——
几分钟之内就会起效，而仍在发育中的年轻\omterm{def:g5:brain-in-command:system}{大脑}受影响
最大。
\end{solution}

\begin{solution}{exo:g5:healthy-choices:6}
身体一旦习惯了某种物质，缺少它就会“抗议”——
于是使用者就被牢牢拴住。开始吸食只是一个简单的选择；
戒掉则是一场年复一年反复进行的抗争：
陷阱的入口很宽，出口很窄，
因此把“永不”的代价定得远高于“以后”。
\end{solution}

\begin{solution}{exo:g5:healthy-choices:7}
事先做出的冷静决定，比在朋友的劝说下临时做出的决定
更容易坚持——别人递过来的那一刻，
恰恰是最不适合开始思考的时候。
\end{solution}

\begin{solution}{exo:g5:healthy-choices:8}
开放性答案。典型的报告是：入睡更快、醒来更轻松，
不看屏幕的那一周早晨精神更好。
\end{solution}

\begin{solution}{exo:g5:healthy-choices:9}
两周的作息时间、饮食和每天的活动都相同——
唯一的变化是：只在第二周睡前一小时关闭屏幕。
测量内容：入睡所需的时间（估计值），醒来时的感觉，
还可以每天记录早晨的情绪。
\end{solution}

\begin{solution}{exo:g5:healthy-choices:10}
以下全部选项。\omterm{prop:g5:healthy-choices:substances}{成瘾}意味着第一支香烟
就是陷阱的入口：所提供的并不只是一口烟，
而是身体对下一支烟需求的开端——
即从容易进入，到艰难脱身。
\end{solution}

\begin{solution}{exo:g5:healthy-choices:11}
气体交换一章解释了清洁空气的界线
和吸烟的警告：\omterm{def:g4:breathing:movements}{呼吸}的作用是让氧气透过娇嫩的表面进入体内，
而烟雾会弄脏这些表面。“心跳加快”的例子解释了运动的界线：
\omterm{def:g4:heart-and-blood:heart}{心脏}和肺是可以训练的——用力就是它们的锻炼。
关于保护\omterm{def:g5:brain-in-command:system}{大脑}的命题解释了\omterm{prop:g1:healthy-habits:needs}{睡眠}的界线
和酒精的警告：指挥中心需要整理归档的时间，
而进入\omterm{def:g5:brain-in-command:system}{大脑}的物质会模糊
那个支配其他一切的循环。
\end{solution}
